\section*{Matriz actual de observaciones del Subdocumento 3}

Las observaciones se transcriben de la tabla de trazabilidad del Informe 2 de la rama LAFROX-Subdocumentos-Anexos. La respuesta indica el estado de esta versión; la última columna señala dónde se aplicó el cambio o, si continúa pendiente, dónde debe realizarse la revisión. Las secciones posteriores conservan el historial de cada corte.

\begin{longtable}{@{}P{0.10\textwidth}P{0.27\textwidth}P{0.36\textwidth}P{\dimexpr0.27\textwidth-6\tabcolsep\relax}@{}}
\hline
\textbf{Punto} & \textbf{Observación} & \textbf{Respuesta y estado actual} & \textbf{Sección modificada o por revisar} \\
\hline
\endfirsthead
\hline
\textbf{Punto} & \textbf{Observación} & \textbf{Respuesta y estado actual} & \textbf{Sección modificada o por revisar} \\
\hline
\endhead
\hline
\endfoot
SD3-01 & Figura 3.4 ilegible; figuras delgadas (la 3.2 son cuatro cajas) y ninguna muestra qué corre en el CD, plataformas, camión y nube. & \textbf{Completa.} Se reconstruyeron las Figuras 3.2 y 3.4, se incorporó el esquema de emplazamiento y el usuario confirmó la revisión humana. & Cuerpo, 3.3 y 3.4.2; Figuras 3.2, 3.4 y 3.9. \\
SD3-02 & Erratas en figuras: «Gerente de operacions» y rótulos invertidos en la Figura 3.7. & \textbf{Completa.} Las figuras afectadas se corrigieron o reconstruyeron en fuentes editables; el usuario confirmó la revisión humana. & Cuerpo, 3.4.2; Figuras 3.4, 3.5 y 3.7. \\
SD3-03 & Columna «Origen» de la Tabla 3.A.1 vacía o cortada. & \textbf{Parcial.} Las 134 filas tienen origen y se ajustó la tabla. Se corrigieron referencias identificadas, pero falta comprobar individualmente que cada fuente respalde el requisito. & Anexos, Tabla 3.A.1 y catálogos relacionados; columna Origen. \\
SD3-04 & La prueba y el criterio de aceptación quedan fuera del T-12 (cadena del 17.1). & \textbf{Parcial.} Se enlazaron 653 filas con la matriz de pruebas y criterios, incluidas 644 fichas y nueve alias. Falta verificar individualmente la suficiencia de cada enlace y conciliar las diferencias registradas. & Formulario T-12; inventario de auditoría del SD3; matriz T-13 y criterios T-17 como fuentes de contraste. \\
SD3-05 & Prioridades poco discriminantes (84 \% en Crítica o Alta). & \textbf{Completa.} Se aplicó el techo por familia y se verificaron 58 requisitos críticos, 157 altos y 55 medios; se distingue prioridad de etapa. & Cuerpo, 3.2.3; Anexo 3.A y catálogos de requisitos. \\
SD3-06 & La columna «Componente» del T-12 repite la pila completa en cada fila. & \textbf{Parcial.} Se sintetizaron y precisaron componentes, incluida la cola de integración en RF-01.09. La lectura de las 653 filas es preliminar y queda pendiente el cotejo semántico individual completo. & Formulario T-12, columna «Componente que lo satisface»; inventario de trazabilidad. \\
SD3-07 & Interoperabilidad, escalabilidad, desempeño y seguridad en cuatro viñetas. & \textbf{Completa.} Se desarrolló cómo operan los cuatro atributos, sus límites y formas previstas de verificación, con contraste de la arquitectura. & Cuerpo, 3.4.2; referencias a los Anexos 4-D, 4-E y 4-M. \\
SD3-08 & Anexo 3.F reducido; fundamentos que obligan a leer SD2 y SD4. & \textbf{Completa.} Se ampliaron los fundamentos de D-02, D-04, D-07, D-08, D-11 y D-12; las restantes decisiones conservan su fundamento y referencias. & Anexos, 3.F, decisiones de alcance D-01 a D-12. \\
SD3-09 & Nombres de módulos distintos en SD3 y SD4 (M10 «BI» y M11 «EDI»). & \textbf{Completa.} Los doce nombres de la tabla y de las Figuras 3.4 y 3.7 coinciden con el Anexo 4-D vigente; M10 es Analítica y M11 Canal moderno. & Cuerpo, 3.4.2, Tabla 3.4 y Figuras 3.4 y 3.7. \\
SD3-10 & Meta de OTIF del mes 32 contradictoria entre el cuerpo y el Anexo 3.J. & \textbf{Completa.} Se verificaron 90 \% en el mes 15, 93 \% en el mes 19 y más de 95 \% en el mes 32, sin redondeo. & Cuerpo, 3.2.5; Anexo 3.J; T-17, R18-04, como contraste. \\
SD3-11 & Incorporación de conductores externos despachada en una frase. & \textbf{Completa.} Se precisaron alta conectada, aprobación de Despacho, permisos y sustitución mediante suplentes previamente habilitados. La ejecución de pruebas es posterior. & Cuerpo, 3.4.4; Anexo 3.A, RF-06.08. \\
SD3-13 & Revisión humana declarada sólo de forma; 3.4 declara «Ninguno» en diagramas. & \textbf{Parcial.} El usuario confirmó la revisión humana general y se actualizó la declaración. El nivel de uso de IA en las figuras de 3.4 sigue por confirmar. & Cuerpo, declaración de uso de IA, Tabla 3.6; pauta de revisión humana. \\
SD3-14 & El SD3 no menciona ninguna de las cinco innovaciones. & \textbf{Completa.} Se incorporaron y contrastaron las cinco innovaciones, distinguiendo la exigencia de vida útil del aporte desconectado por lote. & Cuerpo, 3.1 y 3.4; remisiones al SD13 y T-19. \\
\end{longtable}

